\section{Related work}\label{sec:related}

\textbf{seL4.} seL4 is a capability microkernel with a machine-checked proof that its C implementation refines an abstract specification, and proofs of integrity and confidentiality on top \cite{sel4,sel4info}. \nonos{} shares the capability model and the ring-3 drivers. Its proofs are much weaker in binding: only the capability modules reach extraction, and isolation is proved about models. It adds things seL4 leaves to the system built on it: an admission gate with STARK membership proofs and post-quantum signatures, measured boot checked in both directions, and a desktop and network stack.

\textbf{Fuchsia, Zircon and Starnix.} Zircon is a capability-based kernel in which processes hold handles with rights and drivers run in user space \cite{zircon}. Fuchsia's Starnix runs Linux programs on Zircon by implementing Linux system calls in a user-space component \cite{starnix}. That is the closest relative of the Linux personality, and the shape is the same: the kernel does not speak Linux, a supervisor does. The difference we can point to is the unit of policy. We tie a short signed word to each program's listing and hold the personality's authority to a published capability word. We do not claim that this is better than Starnix's arrangement. We have not compared them.

\textbf{gVisor.} gVisor is an application kernel that intercepts a container's system calls and serves them in user space, to narrow what reaches the host kernel \cite{gvisor}. The personality is the same idea on a smaller host: there is no Linux kernel underneath to protect, because \nonos{} is not Linux. gVisor has far more of Linux behind it than we do. The personality serves 226 of 373 x86\_64 calls by the tree's own coverage tool, and refuses seventeen on purpose.\note{\path{tools/nonos-linux-coverage} prints 226 of 373; \path{docs/userland/linux-personality.md} still says 224.}

\textbf{Capsicum.} Capsicum adds a capability mode to FreeBSD, in which a process loses access to global namespaces and holds rights on file descriptors \cite{capsicum}. The idea that a program gets a short list of rights and nothing ambient is the same. Capsicum asks the program to be written for it, or wrapped. A \nonos{} box is applied to a program that was written for a world of ambient authority, by giving it a world that is mostly empty.

\textbf{Flatpak.} Flatpak sandboxes desktop applications using Linux namespaces and seccomp, with the permissions an application asks for declared in its metadata, and portals for the rest \cite{flatpak}. The declared list and the Marketplace listing play the same role. The difference is the layer: Flatpak's sandbox is a set of kernel features inside a large kernel, and ours is a small kernel plus one supervisor. Flatpak is mature and widely used; the \nonos{} box is new in this release.

\textbf{Qubes OS.} Qubes isolates workloads in Xen virtual machines and routes networking through service VMs, including Whonix for Tor \cite{qubes}. Its isolation boundary is a hypervisor running full guest kernels. \nonos{} runs one small kernel and confines each program by capability, with no guest kernels. The cost is that the compatibility layer is ours to write, and Qubes inherits Linux's.

\textbf{Tails.} Tails is an amnesic live system that routes all traffic through Tor and keeps nothing by default \cite{tails}. \nonos{} shares the amnesic default and the refusal to fall back to the clear network, adds Nym as a default transport, and moves enforcement from a firewall to the route library and, for Linux programs, to the personality.

\textbf{GrapheneOS.} GrapheneOS hardens Android with verified boot, hardware-backed key attestation and a sandboxed app model \cite{graphene}. Its attestation relies on the phone's secure element and the vendor's certificate chain. \nonos{}'s admission proof is checked against a public policy root, and anyone can check it.

\textbf{TPMs and trusted execution.} The TCG specifications define PCR extension and quotes \cite{tcgpc,tcgtpm}. A quote proves what was measured to a verifier who knows the expected values. \nonos{} adds a membership proof so that the expected set is fixed by one root, and an anonymous device proof that hides which member is running. Trusted execution environments such as SGX, TDX and Nitro Enclaves let a remote party verify a measurement signed by hardware keys rooted in the vendor \cite{sgx,tdx,nitro}. They protect a workload from its host but depend on vendor key infrastructure. \nonos{} takes the opposite position: it trusts the CPU and TPM it runs on and makes the software side open and checkable by proofs that need no vendor key.
